
\subsection{Finitary FGLS Monad}
\label{sec:finitary-monad}

A $\phi \in \monFLS\vX\vw$ is finitary when it is generated by its actions on a finite number of \emph{stencils}. 

\begin{definition}[Finitary FGLS Monad]
\[\begin{array}{r@{\;}l@{\,}l@{\,}}
    \monFLSFin         :& \multicolumn{2}{l}{\funC\worldC\setC \to \funC\worldC\setC} \\
    \monFLSFin\vX\vw   \eqdef& \{ \phi \in \monFLS\vX\vw :& \textit{exists a finite family of templates $\Sigma$,} \\
                       &&                                   \textit{a complete cover $\cov$ over $\pair\vw\Sigma$,} \\
                       &&                                   \textit{an $(\vX\anyL\dash \times \theapF)$-valued evaluation $\ev$ over $\Sigma$,} \\
                       &&                                   \textit{such that $\cov,\ev \Vdash \phi$} \}
\end{array}\]
\end{definition}
We give a functor and monad structure to $\monFLSFin$ by showing that the structure of $\monFLS$ lifts through the finitary setting. 

\subsubsection[Action on Morphisms in W]{Action on Morphisms in $\worldC$}

\begin{theorem}
The factorial action of the $\monFLSFin$ monad on morphisms lifts through $\monFLS$, 
that is for every $\sigma : \vw \to \vw'$ we have
\begin{center}
\begin{tikzcd}
    {\monFLSFin\vX\vw}
        \ar[rr, dashed, "\exists\monFLSFin\vX\sigma"] \ar[d, hook]
            && {\monFLSFin\vX\vw'} \ar[d, hook]
    \\
    {\monFLS\vX\vw}
        \ar[rr, "\monFLS\vX\sigma"] 
            && {\monFLS\vX\vw'}
\end{tikzcd}
\end{center}
\end{theorem}
\begin{proof}
Consider a finitary $\phi \in \monFLSFin\vX\vw$, that is, we have a complete cover $\cov$ over $\pair\vw\Sigma$ and an $(\vX\anyL\dash \times \theapF)$-valued evaluation $\ev$ over $\Sigma$ s.t. $\cov,\ev \Vdash \phi$.
The factorial action on $\phi$ is given by precomposition expressed as
\begin{center}
\begin{tikzcd}
    {\homF\worldC{\vw'}{\instU(\dash)} \times \heapF} 
        \ar[rr, "(\dash \comp \sigma) \times \heapF"] \ar[rrr, bend left=15, "\monFLS\vX\sigma(\phi)"]
            && {\homF\worldC\vw{\instU(\dash)} \times \heapF} \ar[r, "\phi"]
                & {\hideMon(\vX\instU \times \heapF)}  \ar[d, "\hmonIso"]
    \\
    {\stencilF{\vw'}\inclITI\dash} \ar[u, "\pair{\anyL\dash}{\mvarM_\dash}"', "\iso"] \ar[rr, "\stencilF\sigma \comp \inclITI\dash"]
                && {\stencilF\vw\inclITI\dash} \ar[u, "\pair{\anyL\dash}{\mvarM_\dash}"', "\iso"] \ar[r, dashed]
                & {\thideMon(\vX\anyL\dash \times \theapF)\inclITI\dash}
\end{tikzcd}
\end{center}

We have to find $\cov^\sigma$, $\ev^\sigma$ s.t. $\cov^\sigma,\ev^\sigma \Vdash \phi \comp (\dash \comp \sigma) \times \heapF$,
So first rebase cover $\cov$ using $\sigma$ (\cref{sec:cov-rebase}),
\begin{center}
\begin{tikzcd}
    {?_\vw}
        \ar[d, "\cov"'] \ar[r, "?_\sigma"]
        & {?_{\vw'}}
            \ar[d, "\cov^\sigma"]
    \\
    {\Sigma} \ar[r, "\rebase^\sigma"' yshift=-2pt]
        & {\Sigma^\sigma} \posquare{ul}
\end{tikzcd}
\end{center}
which gives us a complete cover $\cov^\sigma$.

Define evaluation $\ev^\sigma$ at $\sigma_\vk \in \Sigma^\sigma$,
via $\proj_\vk \comp \rebase^\sigma : \sigma_{\indF{\rebase^\sigma}(\vk)} \to \sigma_\vk$, by 
\[ \proj_\vk \ev^\sigma \mapsto \typeRes{(\vX\anyL\dash \times \theapF)}(\proj_\vk \comp \rebase^\sigma)(\proj_{\indF{\rebase^\sigma}(\vk)} \ev) \]

That $\cov^\sigma,\ev^\sigma \Vdash \phi \comp (\dash \comp \sigma) \times \heapF$, consider $?_{\vw'} \xto\me 1_{\hat\vw}$
\begin{center}
\begin{tikzcd}
    {?_\vw}
        \ar[d, "\cov"'] \ar[r, "?_\sigma"]
        & {?_{\vw'}}
            \ar[d, "\cov^\sigma"] \ar[dr, bend left=30, "\me"]
    \\
    {\Sigma} \ar[r, "\rebase^\sigma"' yshift=-2pt] \ar[rr, dashed, bend right=50, "\widetilde{\me \comp ?_\sigma}"']
        & {\Sigma^\sigma} \ar[r, dashed, "\tilde\me"] \posquare{ul}
            & {1_{\hat\vw}}
\end{tikzcd}
\end{center}
which gives
\begin{align*} 
    \hmonIso(\phi\pair{\anyL{\me \comp ?_\sigma}}{\mvarM_\me})
                    & = \thideMon(\vX\anyL\dash \times \theapF)(\widetilde{\me \comp ?_\sigma})(\proj_{\indF{\widetilde{\me \comp ?_\sigma}}(\star)} \ev) \\
                    & = \thideMon(\vX\anyL\dash \times \theapF)(\tilde\me \comp \rebase^\sigma)(\proj_{\indF{\rebase^\sigma}\indF{\tilde\me}(\star)} \ev) \\
                    & = \thideMon(\vX\anyL\dash \times \theapF)(\tilde\me)(\proj_{\indF{\tilde\me}\star} \ev^\sigma)
\end{align*} 
\end{proof}

\subsubsection[Action on Morphisms in {[W,Set]}]{Action on Morphisms in $\funC\worldC\setC$}
\begin{theorem}
The factorial action of the $\monFLSFin$ monad on morphisms in $\funC\worldC\setC$ lifts through $\monFLS$, 
that is for every $\alpha : \vX \to \vY$ we have
\begin{center}
\begin{tikzcd}
    {\monFLSFin\vX\vw}
        \ar[rr, dashed, "\exists\monFLSFin\alpha_\vw"] \ar[d, hook]
            && {\monFLSFin\vY\vw} \ar[d, hook]
    \\
    {\monFLS\vX\vw}
        \ar[rr, "\monFLS\alpha_\vw"] 
            && {\monFLS\vY\vw}
\end{tikzcd}
\end{center}
\end{theorem}
\begin{proof}
Consider a finitary $\phi \in \monFLSFin\vX\vw$, that is, we have a complete cover $\cov$ over $\pair\vw\Sigma$ and an $(\vX\anyL\dash \times \theapF)$-valued evaluation $\ev$ over $\Sigma$ s.t. $\cov,\ev \Vdash \phi$.

The factorial action on $\phi$ is given by post-composition
\begin{center}
\begin{tikzcd}
    {\homF\worldC\vw{\instU(\dash)} \times \heapF} \ar[r, "\phi"] \ar[rrr, bend left=15, "\monFLS\alpha_\vw(\phi)"]
        & {\hideMon(\vX\instU \times \heapF)} \ar[rr, "\hideMon(\alpha\tinstU \times \heapF)"] \ar[d, "\hmonIso"]
            && {\hideMon(\vY\instU \times \heapF)} \ar[d, "\hmonIso"]
    \\
    {\stencilF\vw\inclITI\dash} \ar[r, dashed] \ar[u, "\pair{\anyL\dash}{\mvarM_\dash}"]
        & {\thideMon(\vX\anyL\dash \times \theapF)\inclITI\dash} \ar[rr, "\thideMon(\alpha\anyL\dash \times \theapF)\comp \inclITI\dash"' yshift=-2pt]
            && {\thideMon(\vY\anyL\dash \times \theapF)\inclITI\dash}
\end{tikzcd}
\end{center}

Define a new evaluation as
\[ \proj_{\sigma_\vi} \ev' \eqdef \typeRes(\alpha\anyL\dash \times \theapF)_{\sigma_\vi}(\proj_{\sigma_\vi} \ev) : \Prod{\sigma_\vi \in \Sigma} \typeRes(\vY\anyL\dash \times \theapF)(\sigma_\vi) \]
Then for $\me :\, ?_\vw \to \inclITI{\hat\vw}$ and induced morphism $\Sigma \xto{\tilde\me} \inclITI{\hat\vw}$.
\begin{align*}
    \hmonIso(\hideMon(\alpha\tinstU \times \heapF)_{\hat\vw}(\phi\pair{\anyL\me}{\mvarM_\me})) & = 
    \thideMon(\alpha\anyL\dash \times \theapF)_{\inclITI{\hat\vw}} (\hmonIso(\phi\pair{\anyL\me}{\mvarM_\me})) \\
        &= \thideMon(\alpha\anyL\dash \times \theapF)_{\inclITI{\hat\vw}}(\thideMon(\vX\anyL\dash \times \theapF)(\tilde\me)(\proj_{\indF{\tilde\me}(\star)} \ev)) \\
        &= \thideMon(\vY\anyL\dash \times \heapF)(\tilde\me)(\thideMon(\alpha\anyL\dash \times \theapF)_{\sigma_\vi}(\proj_{\indF{\tilde\me}(\star)} \ev)) \\
        &= \thideMon(\vY\anyL\dash \times \heapF)(\tilde\me)(\proj_{\indF{\tilde\me}(\star)} \ev') \\
\end{align*}
    and so $\cov, \ev' \Vdash \hideMon(\alpha\tinstU \times \heapF) \comp \phi$.
\end{proof}

\subsubsection{Unit of Monad}

\begin{theorem}
    The unit of the $\monFLSFin$ monad lifts through $\monFLS$, that is we have
\begin{center}
\begin{tikzcd}
    {\vX\vw} \ar[r, dashed, "\eta^\fin_\vw"] \ar[dr, "\eta_\vw"']
        & {\monFLSFin\vX\vw} \ar[d, hook]
    \\
        & {\monFLS\vX\vw}
\end{tikzcd}
\end{center}
\end{theorem}

\begin{proof}
$\eta_\vw(\vx) \in \monFLS\vX\vw$ is given by
\begin{center}
\begin{tikzcd}
    {\homF\worldC\vw{\instU(\dash)} \times \heapF} 
        \ar[rr, "\vX(\dash)(\vx) \times 1"] \ar[rrr, bend left=15, "\eta_\vw(\vx)"]
        && {\vX\instU \times \heapF}
            \ar[r, "\eta^\hideMon"]
            & {\hideMon(\vX\instU \times \heapF)} \ar[d, "\hmonIso"]
    \\
    {\stencilF\vw\inclITI\dash} \ar[u, "\iso"] \ar[rr, "(\vX\anyL\dash(\vx) \times \theapF)\inclITI\dash"]
        && {(\vX\anyL\dash \times \theapF)\inclITI\dash} \ar[r, "\mvarM^{\thideMon} \comp \inclITI\dash"] \ar[u, "\iso"]
            & {\thideMon(\vX\anyL\dash \times \theapF)\inclITI\dash}
\end{tikzcd}
\end{center}

To show that $\eta_\vw(\vx)$ is finitary consider the singleton template $\star \mapsto ?_\vw$,
and the singleton cover over $\pair\vw{\star \mapsto ?_\vw}$, given by the identity stencil $\cov_\star \eqdef\; ?_\vw \xto1 ?_\vw$.
Notice that $\cov$ is an isomorphism, and so it is complete.
Define the $(\vX\anyL\dash \times \theapF)$-valued evaluation at $\star$ to be the resultlet
\[ \proj_\star \ev \eqdef\; \eta_{?_\vw}\pair\vx\emptyset : \typeRes{(\vX\anyL\dash \times \thideMon)}(?_\vw) \]

That $\cov,\ev \Vdash \eta_\vw(\vx)$, consider $?_\vw \xto\me 1_{\hat\vw}$. We have
\begin{align*}
    \hmonIso(\eta_\vw(\vx)\pair{\anyL\me}{\mvarM_\me})
                       & = \eta^\hideMon_{\hat\vw}\pair{\vX\anyL\me(\vx)}{\mvarM_\me} \\
                       & = \eta^\thideMon_{\inclITI{\hat\vw}}\pair{\vX\anyL\me(\vx)}{\mvarM_\me} \\
                       & = \thideMon(\vX\anyL\dash \times \theapF)(\me)(\eta_{?_\vw}\pair\vx\emptyset) \\
                       & = \thideMon(\vX\anyL\dash \times \theapF)(\me)(\proj_\star \ev)
\end{align*}
\end{proof}

\subsubsection{Join of Monad}

\begin{theorem}
    The join of the $\monFLSFin$ monad on morphisms lifts through $\monFLS$.
\begin{center}
\begin{tikzcd}
    {\monFLSFin\monFLSFin\vX\vw}
        \ar[rr, dashed, "\exists \mu^\fin_\vw"] \ar[d, hook]
            && {\monFLSFin\vX\vw} \ar[dd, hook]
    \\
    {\monFLS\monFLSFin\vX\vw}
        \ar[d, "\monFLS\incl_\vw"]
    \\
    {\monFLS\monFLS\vX\vw}
        \ar[rr, "\mu_\vw"] 
            && {\monFLS\vX\vw}
\end{tikzcd}
\end{center}
\end{theorem}
\begin{proof}
    For a finitary $\phi \in \monFLSFin\monFLSFin\vX\vw$, $\mu_\vw(\phi)$ is given by
    \begin{center}
        \begin{tikzcd}
            {\homF\worldC\vw{\instU(\dash)} \times \heapF}
                \ar[r, "\phi"] \ar[rrr, bend left=15, "\mu_\vw(\phi)"]
                & {\hideMon{(\monFLSFin\vX\tinstU \times \heapF)}}
                    \ar[r, "\hideMon(\eval \comp (\vi \times 1))"] \ar[d, "\hmonIso"']
                        & {\Exp\hideMon2{(\vX\tinstU \times \heapF)}}
                            \ar[r, "\mu^\hideMon"]
                            & {\hideMon{(\vX\tinstU \times \heapF)}} \ar[d, "\hmonIso"]
            \\
            {\stencilF\vw\inclITI\dash} \ar[u, "\pair{\anyL\dash}{\mvarM_\dash}", "\iso"'] \ar[r, dashed]
                & {\thideMon(\monFLSFin\vX\anyL\dash \times \theapF)\inclITI\dash} \ar[u, xshift=5pt, "\inv^{\hmonIso}"'] \ar[rr, dashed]
                            && {\thideMon(\vX\anyL\dash \times \theapF)\inclITI\dash}
        \end{tikzcd}
    \end{center}
    which requires
    \begin{enumerate}
        \item A finite family of templates $\Sigma^{\mu\phi}$
        \item A complete cover $\cov^{\mu\phi}$ over $\pair\vw{\Sigma^{\mu\phi}}$
        \item An $(\vX\anyL\dash \times \theapF)$-valued evaluation $\ev^{\mu\phi}$ over $\Sigma^{\mu\phi}$
        \item That $\cov^{\mu\phi},\ev^{\mu\phi} \Vdash \mu_\vw(\phi)$
    \end{enumerate}

    From finitariness of $\phi$, we have a complete cover $\cov$ over $\pair\vw\Sigma$ and $(\monFLSFin\vX\anyL\dash \times \theapF)$-valued evaluation $\ev$ 
    s.t. $\cov, \ev \Vdash \phi$.
    Unpacking this, the resultlet at every $\sigma_\vi \in \Sigma$, $\proj_\vi \ev$ gives
    \begin{itemize}
        \item An allocation morphism $\rho_\vi : \sigma_\vi \ato \sigma_\vi'$.
        \item A finitary computation $\phi_\vi : \monFLSFin\vX\anyL{\sigma_\vi'}$
            \begin{itemize}
                \item A finite family of templates $\Sigma_\vi$
                \item A complete cover $\cov_\vi :\, ?_{\anyL{\sigma_\vi'}} \to \Sigma_\vi$
                \item An evaluation $\ev_\vi : \Prod{\sigma_\vj \in \Sigma_\vi} \typeRes{(\vX\anyL\dash \times \theapF)}(\sigma_\vj)$
            \end{itemize}
        \item A template heap $\mvarM_\vi : \theapF{\sigma_\vi'}$ which is a morphism $?_{\anyL{\sigma_\vi'}} \xto{\mvarM_\vi} \sigma_\vi' \in \tinstC$.
    \end{itemize}
    We lift morphism $\rho_\vi$ to $\tinstC$ by selecting the allocation data using $\dheapF(\concL{\rho_\vi}, 1)(\mvarM_\vi)$,
    and we denote it by $\me_\vi : \sigma_\vi \ato \sigma_\vi' \in \tinstC$.

    \paragraph{Family of templates $\Sigma^{\mu\phi}$}

    We first extract all the ways $\mvarM_\vi$ can unify with cover $\cov_\vi$, that is, we consider the pushout in $\famC\tinstC$
    \begin{center}
        \begin{tikzcd}
            {\Prod{\vi \in \indF\Sigma} ?_{\anyL{\sigma_\vi'}}}
                \ar[rr, "\Prod{\vi \in \indF\Sigma} \cov_\vi"] \ar[dd, "\Prod{\vi \in \indF\Sigma} \mvarM_\vi"'] 
                && {\Prod{\vi \in \indF\Sigma} \Sigma_\vi}
                    \ar[dd, "\Prod{\vi \in \indF\Sigma} \ext^{\cov_\vi}\mvarM_\vi"] 
            \\\\
            {\Prod{\vi \in \indF{\Sigma}} \sigma_\vi'}
                \ar[rr, "\Prod{\vi \in \indF\Sigma}\ext^{\mvarM_\vi}\cov_\vi"']
                && {\Prod{\vi \in \indF\Sigma} \sigma_\vi' \liplus{?_{\anyL{\sigma_\vi'}}} \Sigma_\vi)} \posquare{uull}
        \end{tikzcd}
    \end{center}
    Notice that $\famC\tinstC$ is a coextensive category and thus the pushout above is given as the product of the pushouts at every component.

    Consider a general unifier for a $\mvarM_\vi$, that is consider a $\sigma_\vk \in \sigma_\vi' \liplus{?_{\anyL{\sigma_\vi'}}} \Sigma_\vi$,
    and let $\vj \eqdef \indF{\ext^{\cov_\vi}\mvarM_\vi}(\vk)$ be the index of the stencil in $\cov_\vi$ that matches.
    We have a diagram (connectivity follows from \cref{lem:general-unifier-con-prop})
    \begin{center}
    \begin{tikzcd}
        & {?_{\anyL{\sigma_\vi'}}} \ar[dl, twoheadrightarrow, "\mvarM_\vi"'] \ar[dr, twoheadrightarrow, "\proj_\vj \comp \cov_\vi"]
    \\
    {\sigma_\vi'} \ar[dr, twoheadrightarrow, "\me^\vk_1 \eqdef \proj_\vk \comp \ext^{\mvarM_\vi}\cov_\vi"']
            && {\sigma_\vj} \ar[dl, twoheadrightarrow, "\proj_\vk \comp \ext^{\cov_\vi}\mvarM_\vi \defeq \me^\vk_2"]
    \\
        & {\sigma_\vk}
    \end{tikzcd}
    \end{center}

    Intuitively, $\me^\vk_1$ instantiates the variables $\varL{\sigma_\vi}$ with enough data to match the stencil $\proj_\vj \comp \cov_\vi$.
    But we need to be careful, variables $\varL{\sigma_\vi}$ represent the extensions on the starting heap before the allocation occurs.
    As a result, we need to select only the $\me^\vk_1$ where the instantiated data do not point into the newly allocated data $\mvarM_{\me_\vi}$.
    Semantically, this is expressed as $\me^\vk_1$ should preserve the disconnectivity of $\me_\vi$ (\cref{def:dcon-pres-morph}),
    which is by \cref{lem:dcon-ligu,lem:ligu-dcon} it is equivalent to requiring that
\begin{center}
    \begin{tikzcd}
        & {\sigma_\vi} 
            \ar[dl, twoheadrightarrow, "\con{\me^\vk_1 \comp \me_\vi}"'] \ar[dr, -alloctip, "\me_\vi"]
    \\
    {\sigma_{\con{\me^\vk_1 \comp \me_\vi}}}
        \ar[dr, -rmonotip, "\uncon{\me^\vk_1 \comp \me_\vi}"']
            && {\sigma_\vi'}
                \ar[dl, twoheadrightarrow, "\me^\vk_1"]
    \\
        & {\sigma^\vk_\vi}
    \end{tikzcd}
\end{center}
    is a L.I.G.U.
    When this is the case, the reachable component $\sigma_{\con{\me^\vk_1 \comp \me_\vi}}$ represents only instantiations of the variables $\varL{\sigma_\vi}$
    with data that do not point to the allocated locations.

    Let $\dpres_{\me_\vi}(\sigma_\vi' \pltimes{?_{\anyL{\sigma_\vi'}}} \Sigma_\vi) \subset \sigma_\vi' \pltimes{?_{\anyL{\sigma_\vi'}}} \Sigma_\vi)$
    represent the selection of general unifiers from $\sigma_\vk \in \sigma_\vi' \pltimes{?_{\anyL{\sigma_\vi'}}} \Sigma_\vi$ such that $\me^\vk_1$ preserves the disconnectivity of $\me_\vi$ as described.

    \begin{align*}
        \Sigma 
            \xto{\overbrace{\Prod{\sigma_\vi \in \Sigma} \ipair\vk{\con{\me^\vk_1 \comp \me_\vi}}}^\embed}
        \overbrace{
            \Prod{\sigma_\vi \in \Sigma} \Prod{\sigma_\vk \in \dpres_{\me_\vi}(\sigma_\vi' \pltimes{?_{\anyL{\sigma_\vi'}}} \Sigma_\vi)}\sigma_{\con{\me^\vk_1 \comp \me_\vi}}
        }^{\Sigma^{\mu\phi}} 
    \end{align*}

\paragraph{Cover $\cov^{\mu\phi}$}
    It is enough to show that morphism $\embed : \Sigma \to \Sigma^{\mu\phi}$ is complete, 
    and take the composition $\cov^{\mu\phi} \eqdef\, ?_\vw \xto\cov \Sigma \xto{\embed} \Sigma^{\mu\phi}$ as our cover (complete morphisms are closed under composition).

    That $\embed$ is connected is immediate, so we need to show that is non-overlapping and total.
    For non-overlappingness, it is enough to show that all diagrams of this form satisfy $\vk = \vl$.
    \begin{center}
    \begin{tikzcd}
        & {\sigma_\vi} \ar[dl, "{\con{\me^\vk_1 \comp \me_\vi}}"'] \ar[dr, "\sigma_{\con{\me^\vl_1 \comp \me_1}}"] \ar[dd, "\hat\me"]
    \\
    {\sigma_{\con{\me^\vk_1 \comp \me_1}}} \ar[dr, "\hat\me_\vk"']
            && {\sigma_{\con{\me^\vl_1 \comp \me_\vi}}} \ar[dl, "\hat\me_\vl"]
    \\
        & {\hat\sigma}
    \end{tikzcd}
    \end{center}

    Notice that all the inner squares are L.I.G.U, and by \cref{lem:ligu-pasting,lem:ligu-uniq,lem:ligu-rmono-exist} we have the given isos
    \begin{center}
    \begin{tikzcd}
            && {\sigma_\vi} \ar[dl, twoheadrightarrow, "\con{\me^\vk_1 \comp \me_\vi}"'] \ar[ddll, bend right=40, "\hat\me"'] \ar[ddrr, bend left=40, "\hat\me"] \ar[dr, twoheadrightarrow, "\con{\me^\vl_1 \comp \me_\vi}"] \ar[dd, -alloctip, "\me_\vi"]
    \\
        & {\sigma_{\con{\me^\vk_1 \comp \me_\vi}}} \ar[dl, "\hat\me_\vk"'] \ar[dd, -alloctip, "\uncon{\me^\vk_1 \comp \me_\vi}"']
                && {\sigma_{\con{\me^\vl_1 \comp \me_\vi}}} \ar[dr, "\hat\me_\vl"] \ar[dd, -alloctip, "\uncon{\me^\vl_1 \comp \me_\vi}"]
    \\
    {\hat\sigma} \ar[dd, -alloctip] \ar[dddrr, -alloctip, "\ext^{\hat\me}\me_\vi" near start]
            && {\sigma_\vi'} \ar[ddd, "\ext^{\me_\vi}\hat\me"] \ar[dl, twoheadrightarrow, "\me^\vk_1"'] \ar[dr, twoheadrightarrow, "\me^\vl_1"]
                    && {\hat\sigma} \ar[dd, -alloctip] \ar[dddll, -malloctip, "\ext^{\hat\me}\me_\vi"' near start]
    \\
        & {\sigma^\vk} \ar[dl, crossing over]
                && {\sigma^\vl} \ar[dr, crossing over]
    \\ 
    {\hat\sigma \liplus{\sigma_{\con{\me^\vk_1 \comp \me_\vi}}} \sigma^\vk} \ar[drr, crossing over, "\iso"]
                    &&&& {\hat\sigma \liplus{\sigma_{\con{\me^\vl_1 \comp \me_\vi}}} \sigma^\vl} \ar[dll, crossing over, "\iso"']
    \\
            && {\hat\sigma \liplus{\sigma_\vi} \sigma_\vi'}
    \end{tikzcd}
    \end{center}

    Let $\vj \eqdef \indF{\ext^{\cov_\vi}\mvarM_\vi}(\vk)$ and $\vj' \eqdef \indF{\ext^{\cov_\vi}\mvarM_\vi}(\vl)$ be the indices of their matching stencils.
    We have a diagram
    \begin{center}
    \begin{tikzcd}
        & {?_{\anyL{\sigma_\vi'}}} \ar[dl, "\mvarM_\vi"'] \ar[dr, "\proj_\vj \cov_\vi"'] \ar[drr, "\proj_{\vj'} \cov_\vi"]
    \\
    {\sigma_\vi'} \ar[dr, "\me^\vk_1"'] \ar[drr, "\me^\vl_1" near start]
            && {\sigma_\vj} \ar[dl, crossing over, "\me^\vk_2"' near start]
                & {\sigma_{\vj'}} \ar[dl, "\me^\vl_2"]
    \\
        & {\sigma_\vk} \ar[dr]
            & {\sigma_\vl} \ar[d]
    \\
            && {\hat\sigma \liplus{\sigma_\vi} \sigma_\vi'}
    \end{tikzcd}
    \end{center}
and by $\cov_\vi$ being complete, we have that $\vj = \vj'$, and by \cref{lem:general-unifier-prop}, the general unifiers $\vk, \vl$ must be iso which means they are equal in our setting since $\sigma_\vi' \liplus{?_{\anyL{\sigma_\vi'}}} \Sigma_\vi$ picks representatives for each isomorphism class of general unifiers.

    For totality of $\embed$, we need to show that $\tilde\me$ exists for any $\me$ 
    \begin{tikzcd}
        {\Sigma} \ar[rr, "\me"] \ar[dr, "\embed"']
                && {1_{\hat\vw}}
        \\
            & {\Sigma^{\mu\phi}} \ar[ur, "\exists \tilde\me"']
    \end{tikzcd}

    For $\vi \eqdef \indF\me(\star)$, consider the L.I.G.U. given by \cref{lem:ligu-rmono-exist}
    \begin{tikzcd}
        & {\sigma_\vi} \ar[dl, "\me"'] \ar[dr, -alloctip, "\me_\vi"]
    \\
    {1_{\hat\vw}} \ar[dr, -alloctip, "\ext^\me\me_\vi"']
            && {\sigma_\vi'} \ar[dl, "\ext^{\me_\vi}\me"]
    \\
        & {1_{\hat\vw} \liplus{\sigma_\vi} \sigma_\vi'}
    \end{tikzcd}
    and notice that $1_{\hat\vw} \liplus{\sigma_\vi} \sigma_\vi'$ is a concrete template since it has a variable preserving morphism into it with domain a concrete template.
   
    By totality of $\cov_\vi$, and further by \cref{lem:general-unifier-prop} we have
    \begin{tikzcd} 
        & {?_{\anyL{\sigma_\vi'}}} \ar[dl, "\mvarM_\vi"'] \ar[dr, "\proj_\vj \cov_\vi"] 
    \\
    {\sigma_\vi'} \ar[ddr, "\ext^{\me_\vi}\me"'] \ar[dr, "\me^\vk_1"] 
            && {\sigma_\vj} \ar[dl, "\me^\vk_2"'] \ar[ddl]
    \\
        & {\sigma^\vk} \ar[d, dashed, "\vh"]
    \\
        & {1_{\hat\vw} \liplus{\sigma_\vi} \sigma_\vi}
    \end{tikzcd}
    where $\sigma^\vk$ is its general unifier.
     
    To show that general unifier $\sigma_\vk$ with $\me^\vk_1, \me^\vk_2$ is in $\Sigma^{\mu\phi}$ we only need to show that  
    $\me^\vk_1$ preserves the disconnectivity of $\me_\vi$ which follows from \cref{lem:ligu-dcon,lem:dcon-pres-dclosed}
    \begin{center}
    \begin{tikzcd}
        {\sigma_\vi} \ar[r, -rmonotip, "\me_\vi"] 
            & {\sigma_\vi'} \ar[rr, bend right=20, "\ext^{\me_\vi}\me"'] \ar[r, "\me^\vk_1"] 
                & {\sigma^\vk} \ar[r, dashed, "\vh"] 
                    & {1_{\hat\vw} \liplus{\sigma_\vi} \sigma_\vi'}
    \end{tikzcd}
    \end{center}
    As a result, let the morphism $\tilde\me$ be given by general unifier $\vk$ and the induced morphism by the factorisation system
    \begin{tikzcd}
        {\sigma_\vi} \ar[r, "\me"] \ar[dr, twoheadrightarrow, "\con{\me^\vk_1 \comp \me_\vi}"']
            & {1_{\hat\vw}}
                \ar[r, -alloctip, "\ext^\me\me_\vi"]
                & {1_{\hat\vw} \liplus{\sigma_\vi} \sigma_\vi'}
        \\
            & {\sigma_{\con{\me^\vk_1 \comp \me_\vi}}}
                \ar[u, dashed, "\tilde\me"'] \ar[r, -rmonotip, "\uncon{\me^\vk_1 \comp \me_\vi}"']
                & {\sigma^\vk}
                    \ar[u, "\vh"']
    \end{tikzcd}
      
    \paragraph{Evaluation $\ev^{\mu\phi}$}
    
    We need an evaluation 
    \[ \ev^{\mu\phi} : \Prod{\sigma_{\con{\me^\vk_1 \comp \me_\vi}} \in \Sigma^{\mu\phi}} \typeRes(\vX\anyL\dash \times \theapF)(\sigma_{\con{\me^\vk_1 \comp \me_\vi}}) \]
    We define the resultlet at a $\sigma_{\con{\me^\vk_1 \comp \me_\vi}}$. 
    Let $\vj = \indF{\ext^{\cov_\vi}\mvarM_\vi}(\vk)$ be the index of its matching stencil, and from the diagram
    \begin{tikzcd}
        & {?_{\anyL{\sigma_\vi'}}} \ar[dl, "\mvarM_\vi"'] \ar[dr, "\proj_\vj \cov_\vi"]
    \\
    {\sigma_\vi'} \ar[dr, "\me^\vk_1"']
            && {\sigma_\vj} \ar[dl, "\me^\vk_2"]
    \\
        & {\sigma^\vk}
    \end{tikzcd}
    define the resultlet as
    \[ \proj_{\sigma_{\con{\me^\vk_1 \comp \me_\vi}}} \ev^{\mu\phi} \eqdef 
            \hideRes\relax{\uncon{\me^\vk_1 \comp \me_\vi}} (\typeRes(\vX\anyL\dash \times \theapF)(\me^\vk_2)(\proj_\vj \ev_\vi)) \]

    \paragraph{Realisation $\cov^{\mu\phi},\ev^{\mu\phi} \Vdash \mu_\vw(\phi)$}
    For a $\me : \, ?_\vw \to 1_{\hat\vw}$, we have by totality of $\cov^{\mu\phi}$
    \begin{tikzcd}
        {?_\vw} \ar[r, "\me"] \ar[d, "\cov"']
            & {1_{\hat\vw}}
        \\
        {\Sigma} \ar[ur, "\tilde\me"] \ar[r, "\embed"']
            & {\Sigma^{\mu\phi}} \ar[u, "\me^{\mu\phi}"']
    \end{tikzcd}
    which gives
    \begin{itemize}
        \item An $\vi \eqdef \indF{\tilde\me}(\star)$, and so a $\sigma_\vi \in \Sigma$, and $\tilde\me : \sigma_\vi \to 1_{\hat\vw} \in \tinstC$
        \item A $\vk \eqdef \indF{\me^{\mu\phi}}(\star)$, and a $\me^{\mu\phi} : \sigma_{\con{\me^\vk_1 \comp \me_\vi}} \to \inclITI{\hat\vw}$ with $\tilde\me = \me^{\mu\phi} \comp \embed_\vk$ given by (let $\vj \eqdef \indF{\ext^{\cov_\vi}\mvarM_\vi}(\vk)$)

            \begin{center}
                \begin{tikzcd}
                    & {?_{\anyL{\sigma_\vi'}}} \ar[dl, twoheadrightarrow, "\mvarM_\vi"'] \ar[dr, twoheadrightarrow, "\proj_\vj \, \cov_\vi"]
                \\
                {\sigma_\vi'} \ar[dr, twoheadrightarrow, "\me^\vk_1 \eqdef \proj_\vk \, \ext^{\mvarM_\vi}\cov_\vi"']
                        && {\sigma_\vj} \ar[dl, twoheadrightarrow, "\proj_\vk \, \ext^{\cov_\vi}\mvarM_\vi \defeq \me^\vk_2"]
                \\
                    & {\sigma_\vk}
                \end{tikzcd}
                \begin{tikzcd}
                  {\sigma_\vi} \ar[d, twoheadrightarrow, "\embed_\vk = \con{\me^\vk_1 \comp \me_\vi}"'] \ar[r, -alloctip, "\me_\vi"] \pbsquare{dr}
                      & {\sigma_\vi'} \ar[d, twoheadrightarrow, "\me^\vk_1"]
                  \\
                  {\proj_\vk \Sigma^{\mu\phi} \eqdef \sigma_{\con{\me^\vk_1 \comp \me_\vi}}}
                    \ar[r, -alloctip, "\uncon{\me^\vk_1 \comp \me_\vi}"']
                      & {\sigma^\vk_\vi}
                \end{tikzcd}
            \end{center}
            where the RHS square is a L.I.G.U. by \cref{lem:dcon-ligu}.

        \item Notice that the resultlet $\proj_\vi \ev = \inj_{\sigma_\vi \xato{\rho_\vi} \sigma_\vi'} \pair{\phi_\vi}{\mvarM_\vi}$ extends to
            \[ 
                \typeRes(\monFLSFin\vX\anyL\dash \times \theapF)(\tilde\me)(\proj_\vi \ev) = 
                    \inj_{\inclITI{\hat\vw} \xato{\ext^{\tilde\me}\me_\vi} {\inclITI{\hat\vw} \liplus{\sigma_\vi} \sigma_\vi'}}
                        \pair{\monFLSFin\vX\anyL{\ext^{\me_\vi}{\tilde\me}}(\phi_\vi)}{\theapF(\ext^{\me_\vi}{\tilde\me})(\mvarM_\vi)}
            \]
            using the L.I.G.U. (\cref{lem:ligu-rmono-exist})

            \begin{tikzcd}[sep=small]
                & {\sigma_\vi} \ar[dr, -alloctip, "\me_\vi"]  \ar[dl, "\tilde\me"']
            \\
            {\inclITI{\hat\vw}} \ar[dr, -alloctip, "\ext^{\tilde\me}\me_\vi"']
                    && {\sigma_\vi'} \ar[dl, "\ext^{\me_\vi}{\tilde\me}"]
            \\
                & {\inclITI{\hat\vw} \liplus{\sigma_\vi} \sigma_\vi'}
            \end{tikzcd}
            which further factors as (the outer square is a L.I.S. by \cref{lem:pullback-factoring}, so $\vh$ is induced by \cref{lem:ligu-prop})
            \begin{center}
            \begin{tikzcd}
                {\sigma_\vi} \ar[r, twoheadrightarrow, "\con{\me^\vk_1 \comp \me_\vi}"'] \ar[rr, bend left=25, "\tilde\me"] \ar[d, -alloctip, "\me_\vi"']
                    & {\sigma_{\con{\me^\vk_1 \comp \me_\vi}}} \ar[d, -alloctip, "\uncon{\me^\vk_1 \comp \me_\vi}"] \ar[r, "\me^{\mu\phi}"]
                        & {\inclITI{\hat\vw}} \ar[d, -alloctip, "\ext^{\tilde\me}\me_\vi"]
                \\
                {\sigma_\vi'} \ar[r, "\me^\vk_1"] \ar[rr, bend right=15, "\ext^{\me_\vi}\tilde\me"']
                    & {\sigma_\vk} \ar[r, dashed, "\vh"]
                        & {\sigma_\vi' \liplus{\sigma_\vi} \inclITI{\hat\vw}}
            \end{tikzcd}
            \end{center}
            and notice that the RHS square is itself a L.I.G.U. by \cref{lem:ligu-pasting}.

            Completeness of $\cov_\vi$, gives us
            \[ \hmonIso(\phi_\vi\pair{\anyL{\ext^{\me_\vi}{\tilde\me}}}{\theapF(\ext^{\me_\vi}{\tilde\me})(\mvarM_\vi)}) =
                        \thideMon(\vX\anyL\dash \times \theapF)(\vh)(\typeRes(\vX\anyL\dash \times \theapF)(\me^\vk_2)(\proj_\vj \ev_\vi)) \]
    \end{itemize}

    Lastly we prove $\cov^{\mu\phi},\ev^{\mu\phi} \Vdash \mu_\vw(\phi)$ by
    \begin{align*}
        \hmonIso(\mu_\vw(\phi)\pair{\anyL\me}{\mvarM_\me}) 
            &= \hmonIso(\mu^\hideMon(\hideMon(\eval \comp (\incl \times 1))(\phi\pair{\anyL\me}{\mvarM_\me})))          & \tag{$\cov,\ev \Vdash \phi$} 
\\
            &= \hmonIso(\mu^\hideMon(\hideMon(\eval \comp (\incl \times 1))(
                        \inv^\hmonIso(\thideMon(\monFLSFin\vX\anyL\dash \times \theapF)(\tilde\me)(\proj_\vi \ev)))))
\\
            &= \hmonIso(\mu^\hideMon(\hideMon(\eval \comp (\incl \times 1))(
                    \inv^\hmonIso(\eqclass{\inclITI{\hat\vw} \xato{\ext^{\tilde\me}\me_\vi} {\inclITI{\hat\vw} \liplus{\sigma_\vi} \sigma_\vi'}}
                                          \pair{\monFLSFin\vX\anyL{\ext^{\me_\vi}{\tilde\me}}(\phi_\vi)}{\theapF(\ext^{\me_\vi}{\tilde\me})(\mvarM_\vi)})))) 
\\
            &= \hmonIso(\mu^\hideMon(\eqclass{\inclITI{\hat\vw} \xato{\ext^{\tilde\me}\me_\vi} {\inclITI{\hat\vw} \liplus{\sigma_\vi} \sigma_\vi'}}
                                             \phi_\vi\pair{\anyL{\ext^{\me_\vi}{\tilde\me}}}{\theapF(\ext^{\me_\vi}{\tilde\me})(\mvarM_\vi)})) 
\\
            &= \hmonIso(\hideRes\relax{\ext^{\tilde\me}\me_\vi}\phi_\vi\pair{\anyL{\ext^{\me_\vi}{\tilde\me}}}{\theapF(\ext^{\me_\vi}{\tilde\me})(\mvarM_\vi)})
\\
            &= \hideRes\relax{\ext^{\tilde\me}\me_\vi}(
                    \thideMon(\vX\anyL\dash \times \theapF)(\vh)(\typeRes(\vX\anyL\dash \times \theapF)(\me^\vk_2)(\proj_\vj \ev_\vi)))
                                                                                                                        & \tag{by \cref{lem:thidemon-hide-ext-com}}
\\ 
            &= \thideMon(\vX\anyL\dash \times \theapF)(\me^{\mu\phi})(
                    \hideRes\relax{\uncon{\me^\vk_1 \comp \me_\vi}} (\typeRes(\vX\anyL\dash \times \theapF)(\me^\vk_2)(\proj_\vj \ev_\vi)))
                                                                                                                        & \tag{def $\ev^{\mu\phi}$}
\\ 
            &= \thideMon(\vX\anyL\dash \times \theapF)(\me^{\mu\phi})(\proj_\vk \ev^{\mu\phi})
    \end{align*}
\end{proof}
